\setcounter{Exercise}{0}

\chapter*{Feuille d'exercices : Les conditions}

Vous souhaitez créer un système pour compter les points dans un jeu de dés. Vous lancez 3 dés, vous comptez les points en suivant ces règles :
\begin{itemize}
\item Si vous avez 3 dés différents : vous marquez la somme des dés.
\item Si vous avez 3 dés identiques : vous marquez 30 points
\item Si vous avez 2 dés identiques : vous marquez 20 points
\end{itemize}
\begin{Exercise}[title={Test des dés différents}]
\begin{lstlisting}[style=pythonstyle, language=Python]
# randint(1,6) tire au hasard une valeur entière comprise entre 1 et 6
de1 = randint(1,6)
de2 = randint(1,6)
de3 = randint(1,6)
score = 0
if de1 != de2 and de2 != de3 and de1 != de3 :
    score = de1 + de2 + de3
    print("Vous marquez ",score," points")
else :
    print("Il y a au moins 2 dés identiques")
\end{lstlisting}
\Question Dans ce programme, combien y a-t-il de variables ? Quels sont leurs identifiants ?
\Question Quels opérateurs ont été utilisés ? Que font-ils ?
\Question Imaginons que de1 vaut 1, de2 vaut 4 et de3 vaut 6. Écrire le tableau de variables.
\Question Imaginons que de1 vaut 1, de2 vaut 4 et de3 vaut 4. Écrire le tableau de variables.
\Question Expliquez-en un paragraphe ce que fait le programme.
\end{Exercise}

\newpage
%suivi tableau execution du code
% un si et un sinon, si n 6 et si un trois.
% 
\begin{Exercise}[title={Test des 3 dés identiques}]
\begin{lstlisting}[style=pythonstyle, language=Python]
# randint(1,6) tire au hasard une valeur entière comprise entre 1 et 6
de1 = randint(1,6)
de2 = randint(1,6)
de3 = randint(1,6)
score = 0
if de1 == de2 == de3 :
    score = 30
    print("Vous marquez ",score," points")
else :
    print("Vous n'avez pas eu 3 dés identiques")
\end{lstlisting}
\Question Dans ce programme, combien y a-t-il de variables ? Quels sont leurs identifiants ?
\Question Quel opérateur de comparaison a été utilisé ? Que compare-t-il ?
\Question Imaginons que de1 vaut 1, de2 vaut 4 et de3 vaut 6. Écrire le tableau de variables.
\Question Imaginons que de1 vaut 4, de2 vaut 4 et de3 vaut 4. Écrire le tableau de variables.
\Question Expliquez-en un paragraphe ce que fait le programme.
\end{Exercise}


\begin{Exercise}[title={Test des 2 dés identiques}]
\begin{lstlisting}[style=pythonstyle, language=Python]
# randint(1,6) tire au hasard une valeur entière comprise entre 1 et 6
de1 = randint(1,6)
de2 = randint(1,6)
de3 = randint(1,6)
score = 0
if de1 == de2 :
    score = 20
    print("Vous marquez ",score," points")
elif de1 == de3 :
    score = 20
    print("Vous marquez ",score," points")
elif de2 == de3 :
    score = 20
    print("Vous marquez ",score," points")
else :
    print("Vous n'avez pas eu 2 dés identiques")
\end{lstlisting}
\Question Quel opérateur de comparaison a été utilisé ? Que compare-t-il ?
\Question Imaginons que de1 vaut 3, de2 vaut 4 et de3 vaut 4. Écrire le tableau de variables.
\Question Expliquez-en un paragraphe ce que fait le programme.
\Question Imaginons maintenant que de1 vaut 4, de2 vaut 4 et de3 vaut 4. Écrire le tableau de variables. Que se passe-t-il ?
\end{Exercise}

\newpage
\begin{Exercise}[title={Test des 2 dés identiques}]
Nous voulons maintenant faire un seul script qui affichera le bon score selon le lancer de dé. Voici une proposition de programme :


\begin{lstlisting}[style=pythonstyle, language=Python]
# randint(1,6) tire au hasard une valeur entière comprise entre 1 et 6
de1 = randint(1,6)
de2 = randint(1,6)
de3 = randint(1,6)
score = 0
if de1 != de2 != de3 :
    score = de1+de2+de3
elif de1 == de2 :
    score = 20
elif de1 == de3 :
    score = 20
elif de2 == de3 :
    score = 20
elif de1 == de2 == de3 :
    score = 30
print("Vous marquez ",score," points")
\end{lstlisting}
\Question Imaginons que de1 vaut 3, de2 vaut 6 et de3 vaut 4. Écrivez le tableau de variables.
\Question Imaginons que de1 vaut 3, de2 vaut 4 et de3 vaut 4. Écrivez le tableau de variables.
\Question Imaginons que de1 vaut 4, de2 vaut 4 et de3 vaut 4. Écrivez le tableau de variables.
\Question Est-ce le bon score ? Proposez une modification. \textit{indice : vous n'avez que deux lignes à déplacer !}
\end{Exercise}

\begin{Exercise}[title={Ça se complète}]
\Question Complétez les lignes pointillées \texttt{...} du code ci-dessous demandant un numéro de mois à l'utilisateur et affichant le nombre de jours de ce mois.

On rappelle que Janvier, Mars, Mai, Juillet, Août, Octobre, Décembre ont 31 jours,
Février	a 28/29 jours selon l'année et Avril, Juin, Septembre, Novembre	ont 30 jours.
\begin{lstlisting}[style=pythonstyle, language=Python]
mois = input("Entrez le mois : ")
# On suppose ici que la variable 'mois' est correctement formatée (ex: "Janvier")

if mois == "Janvier" ... mois == "Mars" ... \
   mois == "Mai" ... mois == "Juillet" ... \
   mois == "Août" ... mois == "Octobre" ... \
   mois == "Décembre":
    ...
... mois == "Février":
    jours = 28
elif mois == "Avril" ... mois == "Juin" ... \
     mois == "Septembre" ... mois == "Novembre":
    ...
...:
    jours = "Mois non valide"
print(f"Le mois de {mois} a {jours} jours.")
\end{lstlisting}
\end{Exercise}

\begin{Exercise}[title={À vous de coder}]
\Question Écrivez un programme pour trouver un maximum parmi trois nombres saisis par l'utilisateur.
\end{Exercise}

\begin{Exercise}[title={Résolution d'une équation du second degré}]

L'objectif de cet exercice est d'écrire une fonction Python qui calcule les solutions \textbf{réelles} de l'équation du second degré $ax^2 + bx + c = 0$, en fonction de la valeur du discriminant $\Delta = b^2 - 4ac$.

Selon la nature du discriminant, la formule de recherche des racines est donnée par :
\begin{description}
    \item[Cas 1 : Si le discriminant est positif :] il y a deux racines réelles distinctes données par $\frac{-b+\sqrt{\Delta}}{2a}$ et $\frac{-b-\sqrt{\Delta}}{2a}$.
    \item[Cas 2: Si le discriminant est nul :] il y a exactement une racine réelle donnée par $\frac{-b}{2a}$.
    \item[Cas 3: Si le discriminant est négatif :] il n'y a pas de racines réelles.
\end{description}

\Question
Complétez les lignes pointillées \texttt{...} du code ci-dessous en utilisant les opérateurs de comparaison (comme $<, ==, >$) et l'opérateur \texttt{math.sqrt} (pour la racine carrée). Le module \texttt{math} a déjà été importé.

\begin{lstlisting}[style=pythonstyle, language=Python]
import math

print("Entrer les coefficients")
a = int(input("a : "))
b = int(input("b : "))
c = int(input("c : "))
delta = b**2 - 4*a*c

if ... : # Cas où il n'y a pas de solution réelle
    print("Pas de solution réelle.")
elif ... : # Cas où il y a une seule solution réelle
    x0 = -b / (2*a)
    print(f"Une solution double: x = {x0}")
else: # Cas où il y a deux solutions réelles
    x1 = (-b + math.sqrt(delta)) / (2*a)
    x2 = ...
    print(f"Deux solutions: x1 = {x1} et x2 = {x2}")
\end{lstlisting}

\Question
Écrivez un code Python qui :
\begin{itemize}
    \item Calcule le discriminant $\Delta = b^2 - 4ac$.
    \item Affiche "Delta est positif." et les deux solutions si $\Delta > 0$.
    \item Affiche "Delta est nul." et la solution unique si $\Delta = 0$.
    \item Affiche "Delta est négatif. Les solutions sont complexes (non réelles)." si $\Delta < 0$.
\end{itemize}

\end{Exercise}

% \begin{Exercise}[title={Si vous avez le temps : le Yams ou Yathzee}]

% Le Yam's ou Yahtzee a été créé en 1954 au Canada. Pour jouer au Yams, il faut lancer 5 dés que l'on peut relancer plusieurs fois. Dans cet exercice, nous allons faire un compteur de point pour un lancer de Yams.
% Voici comment se décomptent les point au Yams, d'abord les combinaisons possibles :
% \begin{itemize}
%     \item La combinaison du Brelan : c'est la somme de 3 dés avec des faces identiques.
%     \item La combinaison du Carré : c'est la somme de 4 dés avec des faces identiques
%     \item La combinaison du Full : 3 dés avec des faces identiques, et 2 autres dés avec des faces identiques (25 points).
%     \item La combinaison de la Petite suite : 4 dés dont les chiffres des faces se suivent (20 points).
%     \item La combinaison de la Grande suite : 5 dés dont les faces se précèdent (40 points).
%     \item Et la combinaison du Yam's : lorsque tous les 5 dés présentent des faces identiques (50 points).
% \end{itemize}
% Dans l'absence de combinaison, ont fait la somme des points des dés, dans ce cas :
% un As vaut 1 points, un 2 vaut 2 points, un 3 vaut 3 points, ainsi de suite jusqu'au 6 qui vaut 6 points.
% \Question Écrire un script permettant de calculer le nombre de points d'un lancer au Yams.
% \end{Exercise}

